% ============================================================================
\chapter{Ступень I. Устройство носителя и рождение трёхмерного пространства (\texorpdfstring{$1 + 2^d = d^2$}{1 + 2\^d = d\^2})}
\label{ch:course-stage1-carrier}
% ============================================================================

\section{Урок 1. Задача о самозамкнутом автомате: почему пространство трёхмерно и откуда берётся единица}

Спросим себя: почему пространство вокруг нас имеет именно три измерения ($d = 3$), а не два, не четыре и не одиннадцать?

Представь себе элементарную ячейку решётки в $d$-мерном пространстве $\mathbb{Z}^d$. Чтобы эта ячейка не нуждалась ни в каком «внешнем жёстком диске» или стороннем наблюдателе, её собственный внутренний регистр переходов должен быть \textbf{самозамкнутым}. Откуда вообще берётся бит и что означает самозамкнутость при подсчёте состояний?
\begin{itemize}[noitemsep]
\item \textbf{Неделимость элементарного объёма (откуда берётся $+1$):} частица (возбуждение) занимает весь элементарный объём ячейки целиком и не может занимать его частично. Из невозможности частичного заполнения следует ровно одно неделимое состояние собственной занятости центра $\mathbf{0} = (0, \dots, 0)$.
\item \textbf{Плотное заполнение всех ячеек без пустот (откуда берутся $2^d$ и $1 + 2^d$):} поскольку все элементарные ячейки пространства обязаны примыкать друг к другу вплотную без зазоров, вдоль каждой из $d$ ортогональных осей существуют ровно два плотно сомкнутых противоположных направления $\{-1, +1\}$. Их прямое произведение по $d$ осям образует ровно $2^d$ сомкнутых знаковых октантов $(\pm 1, \dots, \pm 1)$ вокруг центральной ячейки. Вместе с неделимым состоянием собственной занятости центра $(0, \dots, 0)$ получаем ровно \textbf{$|\mathcal{B}(d)| = 1 + 2^d$ неделимых состояний ветвления ячейки}.
\item \textbf{Матрица осевых переходов внутри ячейки (откуда берётся $d \times d = d^2$):} передача состояния от оси к оси внутри клетки без внешней памяти определяется множеством всех упорядоченных пар осей $(i, j) \in \{1, \dots, d\}^2$, мощность которого равна \textbf{$|\mathcal{A}(d)| = d \times d = d^2$} ($d$ продольных пар $i = j$ сохранения оси плюс $d(d-1)$ поворотных пар $i \ne j$ смены оси).
\end{itemize}

Потребуем, чтобы между неделимыми состояниями ветвления плотно заполненной ячейки и парами её осевых переходов существовало взаимно-однозначное соответствие (биекция) --- без информационного дефицита и без мусорной вырожденности:
\begin{equation}
  \label{eq:course-catalan}
  |\mathcal{B}(d)| \;=\; |\mathcal{A}(d)|
  \qquad\Longleftrightarrow\qquad
  1 + 2^d \;=\; d^2.
\end{equation}
Проверим и докажем единственность решения этого уравнения руками:
\begin{itemize}[noitemsep]
\item При $d = 1$: $1 + 2^1 = 3 \ne 1^2 = 1$ (регистра не хватает: на одномерной линии нет поперечных осей для поворота).
\item При $d = 2$: $1 + 2^2 = 5 \ne 2^2 = 4$ (на плоскости снова не хватает одного состояния).
\item \textbf{При $d = 3$: $1 + 2^3 = 1 + 8 = 9 = 3^2$} --- \textbf{точное совпадение!} Ни одного лишнего бита, ни одного недостающего!
\item При $d = 4$: $1 + 2^4 = 17 > 4^2 = 16$, а при переходе $d \mapsto d + 1$ для всех $d \ge 4$ левая часть прирастает на $2^d$, тогда как правая часть прирастает лишь на $2d + 1 < 2^d$. По индукции $1 + 2^d > d^2$ строго для всех $d \ge 4$.
\end{itemize}

\begin{quote}
\small\textbf{Историческая справка: Гипотеза Эжена Каталана (1844) и теорема Преды Михэйлеску (2002).}
В 1844~году бельгийский математик Эжен Шарль Каталан опубликовал в журнале Крелля короткую заметку, в которой высказал гипотезу: среди всех совершенных степеней натуральных чисел $x^a$ и $y^b$ ($x,y,a,b > 1$) существует \emph{только одна-единственная пара последовательных чисел}, отличающихся на единицу: $3^2 - 2^3 = 9 - 8 = 1$. Задача простояла нерешённой $158$~лет, пока в апреле 2002~года румынский математик Преда Михэйлеску не дал её полное доказательство методами круговых полей. Наше уравнение самозамкнутости ячейки $1 + 2^d = d^2 \iff d^2 - 2^d = 1$ --- это диагональное сечение уравнения Каталана ($x = a = d$, $y = 2$), которое выделяет размерность нашего мира $d = 3$ как единственный арифметический резонанс между показательным ветвлением $2^d$ и квадратичной матрицей осей $d^2$.
\end{quote}

Посмотри, что мы только что получили из одной строчки~\eqref{eq:course-catalan}:
\begin{enumerate}[label=\textbf{\arabic*.},noitemsep]
\item Пространство обязано иметь размерность $\mathbf{d = 3}$.
\item Число октантов (и треугольных граней оболочки) равно $\mathbf{|{\triangle}(3)| = 2^3 = 8}$, а число квадратных координатных направлений со знаком равно $\mathbf{|\square(3)| = 2d = 6}$.
\item Разрядность фазового регистра ячейки равна $\mathbf{\mathrm{mod\_bits}(3) = d^2 = 1 + 2^d = 9\text{ бит}}$, а полное число тактов фазового кольца равно:
  \[
    \mathbf{\phaseRingCount(3) \;=\; 2^{\mathrm{mod\_bits}(3)} \;=\; 2^9 \;=\; 512}.
  \]
\item И самое главное: перепишем $1 + 2^d = d^2$ в виде разности:
  \begin{equation}
    \label{eq:course-unit-catalan}
    \mathbf{1 \;=\; d^2 - 2^d \;=\; 3^2 - 2^3 \;=\; 9 - 8 \;=\; \mu_{\mathrm{vac}}(3)}.
  \end{equation}
  Сама математическая единица~$\mathbf{1}$ (вакуумный пол амплитуды $\mu_{\mathrm{vac}}=1$) не постулируется извне --- она есть \textbf{зазор Каталана $d^2 - 2^d$} между квадратом размерности и кубом! В нашем кремниевом чипе \texttt{hardware/cuboctahedron\_carrier\_asic.v} даже константа \texttt{1} нигде не забита числом, а вычисляется в каждом модуле формулой \texttt{localparam integer MU\_VAC = (D * D) - (1 << D);}.
\end{enumerate}

\section{Урок 2. Почему строительный блок Вселенной --- это 12-лучевой кубооктаэдр \texorpdfstring{$O_2$}{O2}}

Теперь посмотрим на элементарный куб Мура $3\times 3\times 3$ ($\{-1,0,+1\}^3$) вокруг узла $(0,0,0)\in\mathbb{Z}^3$. В нём $3^3 - 1 = 26$ соседей, которые под действием поворотов и отражений куба распадаются на три сферические оболочки по квадрату длины вектора $\|e\|_2^2 = e_x^2 + e_y^2 + e_z^2$:
\begin{itemize}[noitemsep]
\item $6$~осевых соседей $O_1$ на расстоянии $\|e\|_2^2 = 1$ (центры $6$ квадратных граней куба $\square$);
\item $12$~рёберных диагональных соседей $O_2$ с одной нулевой и двумя единичными координатами $(\pm 1,\pm 1,0)$, $(\pm 1,0,\pm 1)$, $(0,\pm 1,\pm 1)$ на квадрате расстояния:
  \[
    \|e\|_2^2 \;=\; 1^2 + 1^2 + 0^2 \;=\; \mathbf{2 \;=\; d - 1 \;=\; \Bigl(\frac{\caSignalSpeed}{\speedC}\Bigr)^2};
  \]
\item $8$~угловых соседей $O_3$ на расстоянии $\|e\|_2^2 = 3$ ($8$ вершин куба $\triangle$).
\end{itemize}

\begin{quote}
\small\textbf{Историческая справка: Архимед, Иоганн Кеплер (1611), Леонард Эйлер (1776) и кубооктаэдр.}
Многогранник с $12$ вершинами, $6$ квадратными и $8$ треугольными гранями был описан ещё Архимедом. Своё имя \textbf{кубооктаэдр} он получил потому, что получается пересечением куба ($6$ квадратных граней) и октаэдра ($8$ треугольных граней), если срезать углы куба ровно до середин его $12$ рёбер. В 1611~году Иоганн Кеплер в новогоднем трактате \emph{«О шестиугольных снежинках»} поставил задачу о самой плотной упаковке одинаковых шаров в трёхмерном пространстве и открыл, что вокруг центрального шара вплотную помещаются ровно $12$ таких же шаров, центры которых лежат в вершинах кубооктаэдра (в кристаллографии эта решётка называется \textbf{гранецентрированной кубической, ГЦК / FCC}, а ячейка Вороного--Вигнера--Зейтца вокруг каждого узла --- ромбическим додекаэдром). А в 1776~году Леонард Эйлер доказал, что любой трёхмерный поворот твёрдого тела задаётся \textbf{тремя углами Эйлера} в трёх взаимно перпендикулярных плоскостях $(xy, xz, yz)$. Заметь: число плоскостей вращения $\frac{d(d-1)}{2}$ совпадает с числом пространственных осей~$d$ \emph{только при $d = 3$} ($\frac{3\cdot 2}{2} = 3$), и в каждой из этих трёх плоскостей Эйлера лежат ровно $4$ вершины экваториального квадрата кубооктаэдра ($3 \times 4 = 12$ вершин $O_2(3)$)!
\end{quote}

Докажем, почему именно 12-вершинный \textbf{кубооктаэдр~$\mathbf{O_2}$} является единственной рабочей координационной оболочкой нашего носителя:
\begin{enumerate}[label=\textbf{(\alph*)},noitemsep]
\item \textbf{Единение квадратов и треугольников:} на поверхности кубооктаэдра $O_2$ сходятся ровно $|\square(3)| = 6$ квадратных граней и $|{\triangle}(3)| = 8$ правильных треугольных граней, $12$ вершин ($\geomNeighborTwelve(3) = 2d(d-1) = 12$) и $24$ поверхностных ребра, разбитых на $4$ экваториальных шестиугольника в плоскостях $\{111\}$ решётки ГЦК.
\item \textbf{Изотропия второго и четвёртого моментов:} для любого волнового вектора $\mathbf{k}=(k_x,k_y,k_z)$ сумма квадратов проекций по 12 вершинам кубооктаэдра абсолютно сферически симметрична (в каждую координатную плоскость $xy, xz, yz$ входят по 8 ненулевых компонент $\pm 1$):
  \[
    \sum_{j=1}^{12} (\mathbf{k}\cdot e_j)^2 \;=\; 8\,(k_x^2 + k_y^2 + k_z^2) \;=\; |{\triangle}(3)|\,|\mathbf{k}|^2,
  \]
  а изотропная комбинация четвёртых моментов по оболочкам куба Мура однозначно фиксирует веса $O_2$ относительно $O_1$ и $O_3$.
\item \textbf{Одетое 13-узловое ядро $\mathcal{C}_{13}$ и координационные сферы $S_k = 10k^2 + 2$:}
  вместе с центральным узлом $(0,0,0)$ оболочка $O_2$ образует 13-узловой кластер ($\mathbf{N_\varphi(3) = 1 + 12 = 13}$), а на радиусе $k \ge 1$ шагов по кубооктаэдру число узлов $k$-й координационной сферы выводится по формуле Эйлера $V - E + F = 2$ из роста $6$ квадратных и $8$ треугольных граней:
  \begin{equation}
    \label{eq:course-shells}
    S_k(3) \;=\; 10\,k^2 + 2 \;\;=\;\; \bigl(\geomNeighborTwelve(3) - 2\bigr)k^2 + 2
    \qquad\Longrightarrow\qquad
    S_1 = \mathbf{12},\quad S_2 = \mathbf{42},\quad S_3 = \mathbf{92},\quad S_4 = \mathbf{162},
  \end{equation}
  а полный трёхсферный шар кубооктаэдра содержит ровно:
  \begin{equation}
    \label{eq:course-ball-147}
    B_3(3) \;=\; 1 + S_1 + S_2 + S_3 \;=\; 1 + 12 + 42 + 92 \;=\; \mathbf{147\text{ узлов}}.
  \end{equation}
  Запомни эти четыре числа --- $\mathbf{12, 42, 92, 147}$: мы встретим их и в таблице Менделеева ($Z(\mathrm{Mo})=42$, $Z(\mathrm{U})=92$), и в биологии (147 пар оснований ДНК в нуклеосоме хроматина человека и 42 узла / 147 связей в кортикальной колонке коры мозга!), и в обратном причинном конусе волны!
\item \textbf{Теорема о сокращении масштаба $A = K\cdot dh$, размерность Хаусдорфа $D_H = 3$ вакуумного океана и почему природа полна фракталов ($D_H = \ln N_{\mathrm{sub}}/\ln K < 3$):}
  почему одни и те же числа оболочек кубооктаэдра ($12, 42, 92, 147$) работают и внутри ядра атома, и в молекуле бензола, и в ядре планеты $\mathrm{FeS}$, и в коре головного мозга человека? Потому что при переходе от элементарного шага решётки $dh = 1$ на любой макроэтаж с целочисленным шагом ребра $A = K\cdot dh$ ($K\in\mathbb{N}$) амплитуда масштабируется как $\mu \mapsto K\mu$, билинейные компоненты голономии и знаменателя инерции умножаются на $K^2$, и в законе движения Ньютона--Верле $\Delta^2\varphi = \lfloor N_{\mathrm{torque}}/D(\mathbf{x})\rfloor$ квадрат масштаба $K^2$ \textbf{тождественно сокращается в числителе и знаменателе} ($N_{\mathrm{torque}}^{(K)}\cdot D^{(1)} \equiv N_{\mathrm{torque}}^{(1)}\cdot D^{(K)}$), а квант действия шага $s_{\mathrm{floor}}(K) = (p_0/K)\cdot(K\,dh) = s_0 = 1/2$ и 9-битное кольцо $\mathbb{Z}_{512}$ остаются неизменными!
  При этом формула размерности самоподобия Морана--Хаусдорфа $D_H = \ln N(K)/\ln K$ объясняет сразу два фундаментальных факта о нашем мире (Таблица~\ref{tab:course-hausdorff-fractals}):
  \begin{itemize}[noitemsep]
  \item \emph{Почему само пустое пространство (вакуумный океан ГЦК) трёхмерно и подчиняется закону $1/R^2$ ($D_H = 3$):}
    макро-ячейка Вигнера--Зейтца (ромбический додекаэдр со стороной $A = K\cdot dh$) без единой дыры делится ровно на $N_{\mathrm{pack}}(K) = K^3$ элементарных ячеек ($N_{\mathrm{pack}}(K_1 K_2) = N_{\mathrm{pack}}(K_1)N_{\mathrm{pack}}(K_2)$), поэтому размерность Хаусдорфа самого вакуумного носителя строго целая: $D_H^{(\mathrm{bulk})} = \ln(K^3)/\ln K = \mathbf{3 = d}$;
  \item \emph{Откуда в природе берутся фракталы с дробной размерностью Хаусдорфа ($D_H < 3$):}
    как только включаются правила отбора нашего автомата (световой конус одного такта, порог Гейзенберга $41$, запрет Паули $256$ и ветвление Канала~2), на каждом шаге масштаба~$K$ выживает лишь собственное подмножество $N_{\mathrm{sub}} < K^3$ активных узлов (а $K^3 - N_{\mathrm{sub}} > 0$ узлов вырезаются в фон), порождая на $m$ вложенных этажах настоящие дискретные фракталы с $N_{\mathrm{sub}}^m < K^{3m}$ узлами и дробной размерностью Хаусдорфа $D_H = \ln N_{\mathrm{sub}}/\ln K < 3$!
  \end{itemize}
\end{enumerate}

\begin{table}[htbp]
\centering
\caption{Размерность Хаусдорфа--Морана $D_H = \ln N(K)/\ln K$ сплошного вакуумного носителя ГЦК и четырёх канонических фрактальных каскадов природы.}
\label{tab:course-hausdorff-fractals}
\footnotesize
\begin{tabular}{@{}p{0.25\textwidth}p{0.16\textwidth}p{0.21\textwidth}p{0.31\textwidth}@{}}
\toprule
\textbf{Иерархическая структура} & \textbf{Масштаб $K$ и отбор $N(K)$} & \textbf{Размерность Хаусдорфа $D_H$} & \textbf{Физическое проявление в природе} \\
\midrule
\textbf{1. Вакуумный океан Вороного ГЦК} &
Любое $K\in\mathbb{N}$:\newline $N_{\mathrm{pack}}(K) = K^3$ &
$\dfrac{\ln(K^3)}{\ln K} = \mathbf{3}$ &
Сплошное самоподобное замощение 3D-пространства без дыр; точный закон Гаусса и силы $1/R^2$. \\
\addlinespace[3pt]
\textbf{2. Губка Серпинского--Поста (Канал~2 / 3-SAT)} &
$K = 2$ ($2^3=8$):\newline $N_{\mathrm{sub}} = 7$ ($1$ вырезан) &
$\dfrac{\ln 7}{\ln 2} \approx \mathbf{2{,}8074}$ &
Ветвление обратных конусов прообразов $7^T$ при сбросе вихревого октанта в Канале~2. \\
\addlinespace[3pt]
\textbf{3. Кубооктаэдрическая губка Менгера $\mathcal{C}_{13}$} &
$K = 3$ ($3^3=27$):\newline $N_{\mathrm{sub}} = 13$ ($14$ вырезано) &
$\dfrac{\ln 13}{\ln 3} \approx \mathbf{2{,}3347}$\newline ($O_2$: $\frac{\ln 12}{\ln 3}\approx \mathbf{2{,}2619}$) &
Вложенный светоподобный каркас кубооктаэдра в кубе Мура $3\times 3\times 3$ (вырезанное дополнение $27-13=14=Z(\mathrm{Si})$). \\
\addlinespace[3pt]
\textbf{4. Кластеризация материи на лестнице $\alpha_{\mathrm{fs}}^{-1}$} &
$K = 137$ ($137^3$):\newline $N_{\mathrm{sub}}\in\{13,42,147\}$ &
$\dfrac{\ln N_{\mathrm{sub}}}{\ln 137} \in \mathbf{[0{,}522;\,0{,}999]}$ &
Разреженная иерархия ядер, атомов, кристаллов и галактик в макрообъёме ($f_{\mathrm{occ}}\ll 1$). \\
\addlinespace[3pt]
\textbf{5. Сосудистые и бронхиальные сети жизни} &
$k = d_{\mathrm{in}}(\square) = 4$:\newline $k - \mu_{\mathrm{vac}} = 3$ &
Аллометрия Клейбера:\newline $\alpha = \dfrac{d}{d_{\mathrm{in}}(\square)} = \mathbf{\dfrac{3}{4}}$ &
Самоподобное $C_4$-ветвление кровеносных сосудов, бронхов и деревьев ($B\propto M^{3/4}$). \\
\bottomrule
\end{tabular}
\end{table}

Более того, сами логарифмы $\log_a(b)$, произвольные рациональные степени и корни $a^{p/q}$ и размерности Хаусдорфа--Морана $D_H = \log_K(N_{\mathrm{sub}})$ вычисляются на любом фазовом кольце $\mathbb{Z}_{4Q}$ в целых числах без плавающей точки: разложение $x = 2^k + m$ на характеристику порядка $k = \lfloor\log_2 x\rfloor$ (Канал~2) и фазу мантиссы первого квадранта $\mu = \lfloor m Q / 2^k\rfloor\in[0,Q]$ ($\mathcal{C}_Q(\mu)=Q-\mu,\ \mathcal{S}_Q(\mu)=\mu$, Канал~1) даёт однородный кубический по~$Q$ оператор логарифма
\[
  \widehat{\operatorname{Log2}}_{d,Q}(k,\mu) \;=\; \bigl(d_{\mathrm{in}}(\square)|\Delta(d)|\bigr)Q^2\bigl(k Q + \mathcal{S}_Q(\mu)\bigr) + \mathcal{C}_Q(\mu)\mathcal{S}_Q(\mu)\Bigl(\bigl(|N_{12}(d)|-\mu_{\mathrm{vac}}(d)\bigr)Q + d\bigl(\mathcal{C}_Q(\mu)-\mathcal{S}_Q(\mu)\bigr)\Bigr)
\]
и обратный числитель мантиссы экспоненты $\widehat{\operatorname{Exp2Mant}}_{d,Q}(r) := \bigl(d_{\mathrm{in}}(\square)|\Delta(d)|\bigr)Q^2\mathcal{S}_Q(r) - \mathcal{C}_Q(r)\mathcal{S}_Q(r)\bigl((|N_{12}(d)|-\mu_{\mathrm{vac}}(d))Q - d(\mathcal{C}_Q(r)-\mathcal{S}_Q(r))\bigr)$. При растяжении кольца $Q \mapsto KQ$ и $\mu \mapsto K\mu$ оба оператора умножаются на $K^3$, поэтому в формуле смены основания $\log_a(b) = \widehat{\operatorname{Log2}}_{d,Q}(b)/\widehat{\operatorname{Log2}}_{d,Q}(a)$ масштаб кольца $Q^3$ сокращается тождественно! При $d=3,\ Q=128$ эти формулы дают $\log_2(8)=3$, $\log_3(9)=\log_{10}(100)=2$, $\log_2(3)=203/128$, $\log_2(7)=359/128$, $\log_3(13)=474/203$, $\log_{10}(2)=128/425$, $\sqrt{2}=181/128$, $\sqrt{3}=111/64$, $16^{3/4}=8$ и точно обращают $K^{D_H} = N_{\mathrm{sub}}$: $2^{203/128}=3$, $\operatorname{round}(2^{359/128})=7$ и $3^{474/203}=13$.

\section{Урок 3. Три родные системы счисления кольца \texorpdfstring{$\mathbb{Z}_{512}$}{Z\_512}: 9-битный провод, восьмеричный код \texorpdfstring{$407_8$}{407\_8} и бабочка Грея}

Как записывать числа внутри кольца $\mathbb{Z}_{512}$ так, чтобы вычисления стоили ноль лишних усилий? У нашего 9-битного кольца есть сразу три естественные системы счисления, переходящие друг в друга простой разводкой проводов:

\begin{enumerate}[label=\textbf{\arabic*.},leftmargin=*]
\item \textbf{Родная 9-жильная двоичная шина \texttt{wire [8:0]} (автоматический модуль $512$ без сдвигов и масок):}
  в булевой логике ($\mathbb{Z}_2$) 1-битный провод автоматически складывает числа по модулю~$2$. Точно так же на 9-битной шине \texttt{wire [8:0]} любое сложение и вычитание \textbf{автоматически берёт остаток по модулю $2^9 = 512$} без всяких масок \texttt{\& 511} и без битовых сдвигов! При этом обычный 9-битный сумматор $\texttt{\{carry\_out, sum[8:0]\} = a[8:0] + b[8:0]}$ сам делит результат на два физических канала:
  \begin{itemize}[noitemsep]
  \item \textbf{Канал~1 (9 проводов суммы \texttt{sum[8:0]}):} волновой остаток внутри кольца $\mathbb{Z}_{512}$;
  \item \textbf{Канал~2 (10-й провод переноса \texttt{carry\_out}):} неделимое корпускулярное частное $\lfloor(a+b)/512\rfloor\in\{0,1\}$!
  \end{itemize}
\item \textbf{Восьмеричная триада Вейля--Швингера ($\mathbb{Z}_{512}\cong\mathbb{Z}_8^3$, основание $|{\triangle}(3)| = 8$):}
  поскольку $512 = 2^9 = (2^3)^3 = \mathbf{8^3}$, любое состояние фазы $\varphi\in[0,511]$ записывается ровно \textbf{тремя восьмеричными цифрами}:
  \[
    \varphi \;=\; 64\,d_2 + 8\,d_1 + d_0 \;=\; (d_2\,d_1\,d_0)_8, \qquad d_2,d_1,d_0\in\{0,1,2,3,4,5,6,7\}.
  \]
  Здесь старшая восьмеричная цифра $d_2 = \lfloor\varphi / 64\rfloor\in\mathbb{Z}_8$ --- это квантовый уровень октанта $R$ (в единицах $K_P(3) = 2^{9-3} = \mathbf{64 = 100_8}$), а две младшие цифры $(d_1 d_0)_8\in[0,63]$ --- субквантовая мантисса $q$.
  Посмотри, как в этой родной восьмеричной системе выглядит каноническое начальное слово Вселенной $\mathrm{IV}(3) = 2^{2^3} + (2^3 - 1) = 256 + 7 = 263_{10}$:
  \begin{equation}
    \label{eq:course-iv-octal}
    \mathrm{IV}(3) \;=\; 263_{10} \;=\; 4\cdot 8^2 + 0\cdot 8^1 + 7\cdot 8^0 \;=\; \mathbf{407_8},
  \end{equation}
  где первая цифра $\mathbf{4} = d_{\mathrm{in}}(\square)$ включает полупериод Паули ($4\times 64 = 256$), средняя цифра $\mathbf{0}$ задаёт экваториальный баланс, а младшая цифра $\mathbf{7} = 2^3 - 1$ кодирует $7$ открытых октантов!
\item \textbf{Квадрантно-мантиссовая система Кода Грея ($\mathbb{Z}_{512}\cong\mathbb{Z}_4^{\mathrm{Gray}}\ltimes\mathbb{Z}_{128}$):}
  отделим на 9-битном проводе $u = \varphi[8:0]$ два старших провода квадранта $q_1 = u[8]\in\{0,1\}$ (полукруг $\Phi_{\mathrm{Pauli}}=256$), $q_0 = u[7]\in\{0,1\}$ (четверть круга $Q_\square = 128$) и 7 младших проводов мантиссы $m_7 = u[6:0]\in[0,127]$.
\end{enumerate}

\begin{quote}
\small\textbf{Историческая справка: Рефлексный код Фрэнка Грея (Bell Labs, 1947).}
В обычном двоичном счёте при переходе через границу степени двойки (например, от $127 = 01111111_2$ к $128 = 10000000_2$) одновременно переключаются все восемь бит, что в механических и электронных датчиках угла вызывало ложные промежуточные выбросы. В 1947~году физик лаборатории Bell Labs Фрэнк Грей запатентовал \emph{рефлексный двоичный код} (\emph{Reflected Binary Code}), в котором соседние квадранты отражаются зеркально ($m_7 \leftrightarrow 128 - m_7$), а код квадранта строится через исключающее ИЛИ соседних старших бит: $(g_1, g_0) = (q_1,\ q_1 \oplus q_0)$. В результате на каждом шаге по окружности меняется ровно один бит и исключаются скачки на границах квадрантов.
\end{quote}

